\begin{ex}
Đồng hồ chỉ 8 giờ 15 phút. Kim dài chỉ vào số mấy?
\shortans{3}
\loigiai{
Kim dài chỉ số 3.
}
\end{ex}

\begin{ex}
Đồng hồ chỉ 10 giờ 30 phút. Kim dài chỉ vào số mấy?
\shortans{6}
\loigiai{
Kim dài chỉ số 6.
}
\end{ex}

\begin{ex}
Đồng hồ chỉ 12 giờ đúng. Hai kim đồng hồ cùng chỉ vào số mấy?
\shortans{12}
\loigiai{
Cùng chỉ vào số 12.
}
\end{ex}

\begin{ex}
Lớp học bắt đầu lúc 7 giờ 30 phút và tan học lúc 10 giờ 30 phút. Buổi học kéo dài bao nhiêu giờ?
\shortans{3}
\loigiai{
Kéo dài: $10 - 7 = 3$ giờ.
}
\end{ex}

\begin{ex}
Bé Nam đi ngủ lúc 21 giờ và dậy lúc 6 giờ sáng. Nam đã ngủ bao nhiêu giờ?
\shortans{9}
\loigiai{
Ngủ 9 giờ.
}
\end{ex}

\begin{ex}
Tờ lịch tháng 10 có ngày 20 là Thứ Bảy. Hỏi ngày 21 tháng 10 là ngày gì trong tuần? (Chủ Nhật). Ngày 21 cách ngày 20 bao nhiêu ngày?
\shortans{1}
\loigiai{
Cách nhau 1 ngày.
}
\end{ex}

\begin{ex}
Nếu ngày 5 của tháng là Thứ Hai, thì các ngày Thứ Hai tiếp theo của tháng đó là ngày 12, ngày 19 và ngày bao nhiêu?
\shortans{26}
\loigiai{
Ngày tiếp theo: $19 + 7 = 26$.
}
\end{ex}

\begin{ex}
Tháng đó có tất cả bao nhiêu ngày Thứ Hai (ngày 5, 12, 19, 26)?
\shortans{4}
\loigiai{
Có 4 ngày Thứ Hai.
}
\end{ex}

\begin{ex}
Một tháng có 5 ngày Chủ Nhật, ngày Chủ Nhật đầu tiên là ngày 2. Ngày Chủ Nhật cuối cùng của tháng đó là ngày bao nhiêu?
\shortans{30}
\loigiai{
Các ngày Chủ Nhật: 2, 9, 16, 23, 30. Ngày cuối là ngày 30.
}
\end{ex}

\begin{ex}
Chuyến tàu xuất phát lúc 6 giờ sáng và đến ga lúc 2 giờ chiều (14 giờ). Chuyến tàu đi hết bao nhiêu giờ?
\shortans{8}
\loigiai{
Hết: $14 - 6 = 8$ giờ.
}
\end{ex}

\begin{ex}
Một buổi chiếu phim kéo dài 1 giờ 30 phút. Đổi 1 giờ 30 phút ra phút là bao nhiêu phút?
\shortans{90}
\loigiai{
Ta có: $60 + 30 = 90$ phút.
}
\end{ex}

\begin{ex}
Một tiết học kéo dài 35 phút. Sau 2 tiết học như thế, học sinh đã học tất cả bao nhiêu phút?
\shortans{70}
\loigiai{
Thời gian: $35 + 35 = 70$ phút.
}
\end{ex}

\begin{ex}
Bố đi công tác trong 1 tuần lễ và 3 ngày. Bố đi công tác bao nhiêu ngày?
\shortans{10}
\loigiai{
Đi: $7 + 3 = 10$ ngày.
}
\end{ex}

\begin{ex}
Mẹ đi du lịch trong 2 tuần lễ và 5 ngày. Mẹ đi du lịch tất cả bao nhiêu ngày?
\shortans{19}
\loigiai{
Đi: $14 + 5 = 19$ ngày.
}
\end{ex}

\begin{ex}
Nếu ngày 30 tháng 4 là Thứ Ba, thì ngày 1 tháng 5 tiếp theo là Thứ mấy?
\shortans{4}
\loigiai{
Sau Thứ Ba là Thứ Tư.
}
\end{ex}

\begin{ex}
Nếu ngày 31 tháng 12 là Thứ Sáu, thì ngày 1 tháng 1 năm mới là Thứ mấy?
\shortans{7}
\loigiai{
Sau Thứ Sáu là Thứ Bảy.
}
\end{ex}

\begin{ex}
Đồng hồ điện tử hiển thị 14:30. Đây là mấy giờ rưỡi chiều?
\shortans{2}
\loigiai{
Là 2 giờ rưỡi chiều.
}
\end{ex}

\begin{ex}
Đồng hồ điện tử hiển thị 19:00. Đây là mấy giờ tối?
\shortans{7}
\loigiai{
Là 7 giờ tối.
}
\end{ex}

\begin{ex}
Đồng hồ điện tử hiển thị 22:15. Đây là 10 giờ 15 phút đêm. Số giờ đêm là mấy?
\shortans{10}
\loigiai{
Là 10 giờ đêm.
}
\end{ex}

\begin{ex}
Khoảng thời gian từ 8 giờ sáng đến 12 giờ trưa là bao nhiêu giờ?
\shortans{4}
\loigiai{
Là: $12 - 8 = 4$ giờ.
}
\end{ex}

\begin{ex}
Khoảng thời gian từ 13 giờ đến 17 giờ là bao nhiêu giờ?
\shortans{4}
\loigiai{
Là: $17 - 13 = 4$ giờ.
}
\end{ex}

\begin{ex}
Tháng 8 có 31 ngày. Nếu hôm nay là ngày 25 tháng 8 thì còn bao nhiêu ngày nữa là hết tháng 8?
\shortans{6}
\loigiai{
Còn: $31 - 25 = 6$ ngày.
}
\end{ex}

\begin{ex}
Nếu hôm nay là ngày 10, hôm qua là ngày mấy?
\shortans{9}
\loigiai{
Hôm qua là ngày 9.
}
\end{ex}

\begin{ex}
Nếu hôm nay là ngày 10, ngày kia là ngày mấy?
\shortans{12}
\loigiai{
Ngày kia là ngày 12.
}
\end{ex}

\begin{ex}
Một ngày bác sĩ trực bệnh viện 24 giờ. Bác sĩ trực 2 ngày liên tiếp là bao nhiêu giờ?
\shortans{48}
\loigiai{
Hai ngày là: $24 + 24 = 48$ giờ.
}
\end{ex}
